\documentclass[conference]{IEEEtran}
\IEEEoverridecommandlockouts
% The preceding line is only needed to identify funding in the first footnote. If that is unneeded, please comment it out.
\usepackage{cite}
\usepackage{amsmath,amssymb,amsfonts}
\usepackage{algorithmic}
\usepackage{graphicx}
\usepackage{textcomp}
\usepackage{xcolor}
\def\BibTeX{{\rm B\kern-.05em{\sc i\kern-.025em b}\kern-.08em
    T\kern-.1667em\lower.7ex\hbox{E}\kern-.125emX}}
\begin{document}

\title{
\thanks{Identify applicable funding agency here. If none, delete this.}
}

\author{\IEEEauthorblockN{1\textsuperscript{st} Given Name Surname}
\IEEEauthorblockA{\textit{dept. name of organization (of Aff.)} \\
\textit{name of organization (of Aff.)}\\
City, Country \\
email address or ORCID}
\and
\IEEEauthorblockN{2\textsuperscript{nd} Given Name Surname}
\IEEEauthorblockA{\textit{dept. name of organization (of Aff.)} \\
\textit{name of organization (of Aff.)}\\
City, Country \\
email address or ORCID}}

\maketitle

\begin{abstract}
This survey paper provides a comprehensive overview of the rapidly evolving field of multimodal large language models and their cooperation with agents. The emergence of these models has revolutionized artificial intelligence, enabling agents to interpret and respond to diverse multimodal user queries. Our analysis reveals significant advancements in the architecture and design of multimodal large language models, their applications in vision-language tasks, and the development of training and optimization techniques.

The paper explores the importance of multimodal co-learning and fusion strategies, evaluation metrics and benchmarks, personalization and interaction in multimodal dialogue systems, and challenges and limitations such as data quality, bias, and security concerns. Advanced topics including explainability, transfer learning, and few-shot learning are also discussed, highlighting the complexity and multifaceted nature of this field.

The survey highlights the significance of developing standardized evaluation methodologies and benchmarks to assess the performance of multimodal large language models in collaborative tasks. The future directions and open challenges in this field are identified, with implications for artificial general intelligence, world models, and human-computer interaction.

This paper contributes to the existing literature by providing a thorough analysis of the key themes and findings in multimodal large language models and agent cooperation. The scope of the survey encompasses a wide range of topics, from the foundations of computer vision and natural language processing to advanced techniques such as few-shot learning and explainability. Our study provides valuable insights into the development of complex AI systems and highlights the need for effective coordination and interaction mechanisms between agents. Overall, this survey aims to provide a comprehensive understanding of the current state of multimodal large language models and their cooperation with agents, and to identify future research directions in this field.
\end{abstract}

\begin{IEEEkeywords}
large language models, multimodal learning, natural language processing, machine learning, artificial intelligence
\end{IEEEkeywords}


\section{Introduction to Multimodal Large Language Models}
The emergence of multimodal large language models has revolutionized the field of artificial intelligence, enabling agents to interpret and respond to diverse multimodal user queries \cite{b1}. These models have achieved state-of-the-art performance in powering text-based AI agents, and their extension into the multimodal domain has shown promising results \cite{b2}. The development of collaborative frameworks that integrate multiple large multimodal agents (LMAs) has enhanced collective efficacy, facilitating more effective communication and cooperation between agents \cite{b3}. Furthermore, a comprehensive framework has been proposed to standardize evaluations of LMAs, allowing for more meaningful comparisons and advancements in the field \cite{b4}.

Multimodal understanding and generation are key aspects of these models, highlighting the importance of combining multiple data types to enable a more comprehensive understanding of diverse data \cite{b5}. Large language models (LLMs) serve as a foundation for multimodal agents, and their use has become a common theme across various studies \cite{b6}. Collaboration and communication between agents, including agent-agent and agent-human interactions, are critical aspects of multimodal agents, emphasizing the need for effective communication protocols \cite{b7}. The evaluation and benchmarking of LMAs are also crucial, with a need for standardized methodologies and frameworks to facilitate meaningful comparisons \cite{b8}.

While some studies focus primarily on multimodal understanding, others emphasize multimodal generation and editing, highlighting the diversity of approaches in this field \cite{b9}. The use of proprietary models, such as GPT-4o, has been mentioned in some studies, raising concerns about potential limitations and biases associated with these models \cite{b10}. Additionally, there is a distinction between agent-centric and human-centric approaches, emphasizing the importance of considering both aspects in multimodal agent development \cite{b11}.

The technical details and methodologies employed in the development of multimodal large language models are varied, including different architectures and training strategies \cite{b12}. Various datasets have been used for evaluation and benchmarking purposes, highlighting the need for high-quality, diverse data to train effective multimodal agents \cite{b13}. Techniques such as Multimodal ICL (M-ICL), Multimodal CoT (M-CoT), and LLM-Aided Visual Reasoning (LAVR) have been introduced, demonstrating the innovative approaches being explored in this field \cite{b14}.

Despite the advancements in multimodal large language models, several limitations and challenges remain. The need for standardized evaluation methodologies and benchmarking frameworks is a significant challenge, hindering meaningful comparisons and progress in the field \cite{b15}. The limited availability of high-quality, diverse datasets for multimodal agent training is another challenge, emphasizing the need for more extensive data collection and curation efforts \cite{b16}. Furthermore, the complexity and scalability of multimodal agents are significant challenges, requiring innovative solutions to facilitate more efficient and effective development \cite{b17}. Finally, the need for more explainable and transparent multimodal agents is essential to facilitate trust and adoption, highlighting the importance of developing models that can provide clear insights into their decision-making processes \cite{b18}.

In conclusion, the development of multimodal large language models has made significant progress in recent years, with various studies exploring different approaches and techniques. However, several challenges and limitations remain, emphasizing the need for continued research and innovation in this field. The cooperation between multimodal large language models and agents is a critical aspect of this research, highlighting the potential benefits of integrating these technologies to facilitate more effective communication and collaboration \cite{b19}. As the field continues to evolve, it is essential to address the challenges and limitations associated with multimodal large language models, ultimately leading to the development of more effective and transparent models that can facilitate meaningful interactions between humans and agents.

\section{Background and Foundations: Computer Vision and Natural Language Processing}
The intersection of computer vision and natural language processing has given rise to multimodal large language models (MLLMs) and agents, enabling them to interpret and respond to diverse user queries. A comprehensive analysis of relevant papers reveals a complex landscape of research focused on the foundations and applications of MLLMs. For instance, the work presented in \cite{b1} provides a thorough survey of multimodal foundation models, highlighting their evolution from specialist models to general-purpose assistants. This transition underscores the growing importance of multimodality in AI systems, allowing them to process and understand multiple forms of data, including images, text, and speech.

The architectures, applications, and impact of MLLMs on AI and generative models are explored in \cite{b2}, offering a balanced perspective on opportunities and challenges. This study emphasizes the potential of MLLMs to revolutionize vision-language tasks, such as image captioning, visual question answering, and text-to-image synthesis. However, it also notes the difficulties in evaluating these models, particularly in terms of their ability to generalize across different domains and tasks. The analysis presented in \cite{b3} introduces a systematic framework to quantify capability disparities between different modalities in MLLMs, revealing that models like GPT-4V tend to perform consistently across modalities for simple tasks but struggle with more challenging ones.

The development of open-source datasets, frameworks, and benchmarks is crucial for training and evaluating MLLMs. The work presented in \cite{b4} introduces LAMM, a language-assisted multi-modal instruction-tuning dataset, framework, and benchmark, enabling seamless human-AI interaction. This contribution highlights the need for standardized evaluation methods, as well as the importance of multimodality in facilitating effective communication between humans and AI agents. A systematic review of large multimodal agents is conducted in \cite{b5}, categorizing the current body of research into four distinct types and proposing a comprehensive framework to standardize evaluations.

A common thread throughout these papers is the emphasis on multimodality, MLLMs, and the need for standardized evaluation methodologies. The importance of human-AI interaction is also a recurring theme, with studies like \cite{b4} and \cite{b5} emphasizing its significance in enabling AI agents to interpret and respond to diverse user queries. However, contrasting viewpoints emerge regarding the transition from specialist models to general-purpose assistants, as well as the evaluation frameworks proposed by different researchers. For example, \cite{b1} highlights the benefits of general-purpose models, while other papers focus on specific applications or evaluations.

The technical details and methodologies presented in these studies provide valuable insights into the construction of multi-modal instruction tuning datasets and benchmarks, model training, and evaluation metrics. The work presented in \cite{b4} outlines a methodology for constructing such datasets and benchmarks, while also providing a primary but potential MLLM training framework optimized for modality extension. Papers like \cite{b5} and \cite{b2} discuss various evaluation metrics and methodologies, highlighting the need for standardized approaches to assess MLLMs and agents.

Despite the progress made in this field, several limitations and challenges remain. The use of diverse evaluation methods across existing studies hinders effective comparison among different MLLMs, as noted in \cite{b5}. Furthermore, the scalability and extension of MLLM research to diverse domains, tasks, and modalities are essential for rapid progress, as emphasized in \cite{b4}. The capability disparities between different modalities, revealed in \cite{b3}, also indicate a need for further research on this topic. Addressing these challenges will be crucial for the development of more effective and generalizable MLLMs, ultimately enabling more seamless human-AI interaction and expanding the potential applications of these technologies.

\section{Architecture and Design of Multimodal Large Language Models}
The architecture and design of multimodal large language models (MLLMs) have undergone significant transformations in recent years, driven by the need for more sophisticated and interactive AI systems. A key finding in this area is the development of large multimodal agents (LMAs) that can interpret and respond to diverse multimodal user queries \cite{b1}. This has been achieved through the introduction of novel benchmarks, such as COMMA, which evaluates the collaborative performance of multimodal multi-agent systems through language communication \cite{b2}. The application of MLLMs to various vision-language tasks, including visual storytelling and enhanced accessibility, has also shown promising results \cite{b3}.

The importance of developing multimodal agents that can effectively communicate with humans and other agents is a common theme across the literature. This highlights the need for standardized evaluation methodologies to compare the performance of different LMAs \cite{b1}, \cite{b2}. Furthermore, the potential applications of MLLMs in various fields, including computer vision, natural language processing, and game playing, are emphasized in several studies \cite{b3}, \cite{b4}. However, the challenges associated with scalability, robustness, and cross-modal learning in MLLMs remain a concern \cite{b3}, \cite{b5}.

The technical approaches employed in the development of MLLMs vary widely, ranging from architectural components for developing LMAs \cite{b1} to training methods such as cross-modal learning and self-supervised learning \cite{b3}. Evaluation methodologies also differ, with some studies introducing novel benchmarks \cite{b2} while others focus on evaluating MLLMs in specific tasks \cite{b3}. The use of LLM-based and CLIP/T5-based methods for multimodal generation and editing has also been explored \cite{b6}.

Despite the progress made, several limitations and challenges remain. Scalability and robustness issues in cross-modal learning are a significant concern \cite{b3}, while the need for large-scale datasets for training and evaluating MLLMs is also highlighted \cite{b1}. The development of effective evaluation methodologies for assessing the performance of LMAs is another challenge \cite{b2}, as is the importance of addressing safety concerns in multimodal generation and editing, such as avoiding biased or toxic content \cite{b6}.

The implications of these developments are significant, with potential applications in a wide range of fields. The ability of MLLMs to effectively communicate with humans and other agents has far-reaching consequences for areas such as human-computer interaction, computer vision, and natural language processing. Furthermore, the connection between MLLMs and multimodal generation and editing highlights the importance of considering the safety and ethical implications of these technologies \cite{b6}. As research in this area continues to evolve, it is likely that we will see significant advancements in the development of more sophisticated and interactive AI systems, with potential applications in a wide range of fields.

\section{Applications of Multimodal Large Language Models: Vision-Language Tasks}
The applications of multimodal large language models (MLLMs) in vision-language tasks have witnessed significant growth, with a multitude of research papers highlighting their potential in image captioning, object detection, visual storytelling, and other related areas \cite{b1}. A key finding across these studies is the emergence of surprising capabilities in MLLMs, such as generating stories based on images and performing OCR-free math reasoning, which suggests a promising path towards achieving artificial general intelligence \cite{b2}. These advancements underscore the importance of understanding cross-modal interactions between different modalities in MLLMs, with proposals for methods to quantify capability disparities between modalities being an active area of research \cite{b3}.

The dominance of vision-language tasks as a primary application area for MLLMs is evident, reflecting a growing interest in developing general-purpose assistants that can seamlessly integrate multiple modalities, including text, images, video, and audio \cite{b4}. Researchers are exploring various architectures, training methods, and evaluation metrics to improve the performance of MLLMs in these tasks. For instance, instruction-tuning in large language models has been proposed as a strategy for building multimodal models akin to GPT-4, showcasing the versatility and potential of such approaches \cite{b5}. However, this also introduces considerations regarding the ethical implications of developing and deploying MLLMs, emphasizing the need for responsible AI development and a nuanced understanding of future directions \cite{b6}.

While some studies focus on developing specialized models tailored to specific vision-language tasks, others advocate for a more general-purpose approach to building MLLMs capable of handling a wide range of tasks \cite{b7}. This dichotomy reflects differing viewpoints within the research community, with some emphasizing the need for a deeper understanding of cross-modal interactions and others prioritizing improvements in overall performance metrics \cite{b8}. Technical details vary widely, with MLLMs often involving the integration of multiple modalities through unified vision models or end-to-end training, and training methods encompassing pre-training on large datasets, fine-tuning for specific tasks, and instruction-tuning \cite{b9}.

Despite these advancements, several challenges persist, including scalability and robustness, which are critical for developing MLLMs that can effectively handle large-scale datasets and diverse tasks \cite{b10}. Cross-modal learning and understanding the intricate interactions between different modalities remain open research questions, necessitating further investigation \cite{b11}. Additionally, ethical considerations such as bias, fairness, and transparency are paramount when developing and deploying MLLMs in real-world applications, given the potential for hallucinations or errors in multimodal outputs \cite{b12}. The rapid evolution of MLLMs underscores the importance of addressing these challenges to fully harness their potential in vision-language tasks and beyond.

In conclusion, the applications of MLLMs in vision-language tasks represent a vibrant and rapidly evolving field, marked by significant achievements and open research questions. As researchers continue to explore the capabilities and limitations of MLLMs, it is crucial to consider the broader implications of these developments, including their potential for enhancing human-agent cooperation and contributing to the advancement of artificial intelligence as a whole \cite{b13}. By addressing the challenges associated with MLLMs and pursuing a deeper understanding of cross-modal interactions and ethical considerations, the research community can unlock the full potential of these models, paving the way for innovative applications in vision-language tasks and other areas.

\section{Training and Optimization Techniques for Multimodal Models}
Training and optimization techniques for multimodal models have become a crucial aspect of research in recent years, with a focus on developing large multimodal agents (LMAs) that can interpret and respond to diverse multimodal user queries \cite{b1}. The importance of data processing techniques in modern multimodal model training cannot be overstated, with diffusion models and multimodal large language models (MLLMs) playing a significant role \cite{b2}. Researchers have emphasized the need for effective data processing techniques to improve multimodal model performance, highlighting the importance of incorporating multiple modalities to enhance generation ability and interaction with the world \cite{b3}.

The development of a unified multimodal model remains a significant challenge, requiring integration of various techniques and technologies \cite{b4}. To address this challenge, researchers have proposed several approaches, including multimodal chain of thought (M-COT), multimodal instruction tuning (M-IT), and multimodal in-context learning (M-ICL) \cite{b5}. These techniques aim to achieve artificial general intelligence and world models by enabling multimodal models to reason and make decisions more effectively. Furthermore, tool-augmented multimodal agents have been proposed as a means to leverage existing generative models for human-computer interaction, highlighting the potential of multimodal learning in real-world applications \cite{b6}.

Despite the advancements in training and optimization techniques, several challenges persist, including the lack of large-scale datasets for training and testing multimodal models \cite{b7}. The evaluation of different multimodal models is also hindered by diverse evaluation methodologies, requiring the establishment of standardized frameworks \cite{b8}. Moreover, ensuring the safety and robustness of multimodal generation and editing models is a critical challenge, particularly in applications where factuality and reasoning are essential \cite{b9}.

To overcome these challenges, researchers have emphasized the need for designing effective evaluation methodologies and frameworks to compare different multimodal models in a standardized manner \cite{b10}. Additionally, creating tool-augmented multimodal agents that can leverage existing generative models for human-computer interaction has been proposed as a potential solution \cite{b11}. The integration of external rule systems to improve reasoning and decision-making in multimodal models is also an active area of research, with several studies exploring the use of logical rules and constraints to enhance model performance \cite{b12}.

The implications of these developments are significant, highlighting the need for continued innovation in training and optimization techniques for multimodal models. By addressing the challenges associated with data processing, evaluation methodologies, and safety, researchers can develop more sophisticated multimodal models that can unlock the full potential of multimodal learning and achieve artificial general intelligence \cite{b13}. Ultimately, the development of effective training and optimization techniques will be crucial in enabling multimodal large language models and agents to cooperate together seamlessly, leading to significant advancements in human-computer interaction and artificial intelligence \cite{b14}.

\section{Multimodal Co-learning and Fusion Strategies}
Multimodal co-learning and fusion strategies have emerged as crucial components in the development of large language models and agents that can effectively cooperate together. Recent studies, such as those presented in \cite{b1} and \cite{b2}, have highlighted the significance of evaluating collaborative performance in multimodal multi-agent systems through language communication, underscoring weaknesses in state-of-the-art models in agent-agent collaborations. The introduction of novel benchmarks like COMMA \cite{b3} has further emphasized the need for comprehensive evaluation frameworks that can assess the performance of multimodal models and agents in cooperative and heterogeneous settings.

A key theme that emerges from the analysis of relevant papers is the importance of addressing the challenges posed by multimodality, including modality heterogeneity, connections, interactions, and the impact of missing or noisy data \cite{b4}. The potential benefits of co-learning strategies that enable knowledge transfer between different modalities are also emphasized, particularly in addressing scenarios where one modality may have limited resources or data \cite{b5}. For instance, techniques for transferring knowledge between modalities can improve model performance in resource-constrained scenarios, as discussed in \cite{b6}.

The development of multimodal agents that can understand and interact with their environment through multiple modalities, such as text and vision, is also a notable trend \cite{b7}. These agents can enhance their decision-making and reasoning abilities by leveraging the strengths of different modalities. However, this also introduces new challenges, such as handling missing or noisy modalities, understanding complex interactions between modalities, and ensuring scalability and generalizability \cite{b8}.

Despite these challenges, multimodal co-learning and fusion strategies offer significant opportunities for advancing research in this area. The use of benchmarks like CH-MARL \cite{b9} can facilitate the development of more effective cooperative, heterogeneous multi-agent learning frameworks. Furthermore, the emphasis on comprehensive evaluation frameworks and the importance of standardization can help to establish a common foundation for comparing different models and approaches \cite{b10}.

In conclusion, multimodal co-learning and fusion strategies are critical components in the development of large language models and agents that can cooperate effectively. By addressing the challenges posed by multimodality and leveraging the potential benefits of co-learning strategies, researchers can create more effective and scalable multimodal models and agents. The implications of this research are significant, with potential applications in areas such as human-computer interaction, robotics, and natural language processing \cite{b11}. As the field continues to evolve, it is essential to prioritize standardization, handle missing or noisy data, and develop more sophisticated models that can capture complex interactions between modalities. Ultimately, the development of multimodal co-learning and fusion strategies will be crucial for realizing the full potential of large language models and agents in cooperative and heterogeneous settings \cite{b12}.

\section{Evaluation Metrics and Benchmarks for Multimodal Large Language Models}
Evaluation Metrics and Benchmarks for Multimodal Large Language Models

The development of multimodal large language models (MLLMs) has led to a growing need for effective evaluation metrics and benchmarks to assess their performance. Recent studies, including COMMA \cite{b1}, MM-BigBench \cite{b2}, ChEF \cite{b3}, and AgentBoard \cite{b4}, have made significant contributions to this area by introducing novel benchmarks and frameworks for evaluating MLLMs. These contributions highlight the weaknesses in state-of-the-art models, such as the lack of standardized evaluation frameworks, and emphasize the importance of establishing comprehensive assessment tools.

A key finding across these studies is the need for standardized evaluation frameworks that facilitate fair comparisons and meaningful evaluations of MLLMs. For instance, ChEF \cite{b3} proposes a Comprehensive Evaluation Framework for standardized assessment of MLLMs, enabling holistic profiling and fair comparison of different models. Similarly, MM-BigBench \cite{b2} presents a comprehensive assessment framework for evaluating MLLMs on multimodal content comprehension tasks, incorporating diverse metrics to achieve a more holistic evaluation. These frameworks address the challenges in evaluating MLLMs and provide more effective benchmarks for assessing their performance.

The importance of multimodal content comprehension is another common theme emerging from these studies. Papers like MM-BigBench \cite{b2} and ChEF \cite{b3} highlight the need to evaluate MLLMs on tasks that require a deeper understanding of multimodal contexts, such as image-text retrieval and visual question answering. Additionally, COMMA \cite{b1} and AgentBoard \cite{b4} focus on evaluating the collaborative performance of MLLMs in multi-agent settings, emphasizing the importance of language communication and cooperation.

While these studies share common themes and goals, some contrasting viewpoints and approaches are evident. For example, COMMA \cite{b1} and MM-BigBench \cite{b2} focus on developing specific benchmarks for evaluating MLLMs, whereas ChEF \cite{b3} and AgentBoard \cite{b4} propose more general frameworks for standardized evaluation. Furthermore, the emphasis on language communication versus multimodal content comprehension varies across studies, reflecting the diversity of perspectives and methodologies in the field.

The technical details and methodologies employed in these studies demonstrate the complexity and nuance involved in evaluating MLLMs. For instance, MM-BigBench \cite{b2} and ChEF \cite{b3} involve the creation of large-scale datasets for evaluating MLLMs on multimodal content comprehension tasks, while AgentBoard \cite{b4} provides a fine-grained progress rate metric for analyzing multi-turn interactions between LLM agents. These technical details highlight the need for careful consideration of evaluation design and implementation.

Despite the significant contributions made by these studies, several limitations and challenges remain. The absence of standardized evaluation frameworks makes it difficult to compare and contrast different MLLMs \cite{b1}. Additionally, evaluating MLLMs on large-scale datasets and ensuring generalizability across diverse scenarios remain significant challenges \cite{b2}. AgentBoard \cite{b4} also highlights the need for more interpretable evaluation metrics and tools to provide deeper insights into LLM agent behaviors.

In conclusion, the development of effective evaluation metrics and benchmarks for MLLMs is an active area of research. The studies discussed in this section highlight the need for standardized evaluation frameworks, the importance of multimodal content comprehension, and the challenges involved in evaluating MLLMs. As the field continues to evolve, it is essential to address these challenges and develop more comprehensive and standardized approaches to evaluating MLLMs. By doing so, we can facilitate meaningful comparisons and advancements in the development of MLLMs, ultimately leading to improved performance and applications in various domains \cite{b5}.

\section{Personalization and Interaction in Multimodal Dialogue Systems}
Personalization and interaction are crucial aspects of multimodal dialogue systems, enabling effective communication between humans and agents. Recent research has made significant contributions to our understanding of these concepts, highlighting the importance of developing unified languages for communicative acts between agents \cite{b1}. For instance, the introduction of a multimodal meaning representation framework for generic dialogue systems has paved the way for designing multi-agent architectures that facilitate seamless interaction \cite{b2}. Furthermore, persona-conditioned large language models (LLMs) have been shown to engage in open and naturalistic dialogues when personality consistency and linguistic alignment are ensured \cite{b3}.

The development of novel benchmarks, such as COMMA, has also enabled the evaluation of collaborative performance in multimodal multi-agent systems through language communication \cite{b4}. Customizing agent personas has been found to lead to improved interaction quality, diversity, and dynamics, as demonstrated by the CloChat platform, which supports easy and accurate customization of agent personas in LLMs \cite{b5}. A comprehensive survey of large multimodal agents (LMAs) has underscored the need for standardization in evaluation methodologies to facilitate meaningful comparisons between different approaches \cite{b6}.

A common theme emerging from these studies is the importance of frameworks and architectures that support task-independent dialogue managers and enable effective communication between agents. Personalization and customization in LLMs are also recognized as essential for improving interaction quality and user experience. The growing interest in multimodal interaction and the extension of LLMs into the multimodal domain is another notable trend \cite{b7}. However, challenges persist in evaluating and comparing different approaches, highlighting the need for standardized evaluation methodologies \cite{b8}.

Different studies have adopted contrasting viewpoints or approaches, such as the focus on developing a unified language for communicative acts versus the emphasis on customizing agent personas \cite{b9}. The use of different evaluation methodologies, such as COMMA and Large Multimodal Agents: A Survey, has also been noted \cite{b10}. Technical details and methodologies employed in these studies include the development of meta-models, prompting and sampling algorithms, and novel benchmarks \cite{b11}. Despite these advancements, limitations and challenges remain, including the need for more research on task-independent dialogue managers and standardized evaluation methodologies \cite{b12}.

The implications of these findings are significant, suggesting that personalized and interactive multimodal dialogue systems have the potential to revolutionize human-agent communication. By developing more effective frameworks and architectures, customizing agent personas, and standardizing evaluation methodologies, researchers can create more naturalistic and engaging interactions between humans and agents. Future studies should investigate the impact of customizing agent personas on interaction quality and user experience, as well as explore the development of more advanced meta-models and evaluation methodologies \cite{b13}. Ultimately, the integration of multimodal large language models and agents has the potential to enable more effective and personalized communication, with far-reaching implications for fields such as human-computer interaction, natural language processing, and artificial intelligence.

\section{Challenges and Limitations: Data Quality, Bias, and Security Concerns}
The integration of multimodal large language models (MLLMs) and agents is a rapidly evolving field, with significant potential for enhancing collaborative tasks, situational safety assessments, and real-world fact-checking applications. However, as highlighted by recent studies, including \cite{b1}, \cite{b2}, and \cite{b3}, several challenges and limitations must be addressed to ensure the reliable and trustworthy deployment of these models. A primary concern is the quality and availability of high-quality data, which is essential for improving MLLM performance \cite{b4}. The Synergy between Data and Multi-Modal Large Language Models survey emphasizes the interconnected development of models and data, underscoring the need for co-development strategies to enhance MLLM capabilities \cite{b5}.

The lack of standardized evaluation methods is another significant challenge, hindering effective comparison and improvement of multimodal models \cite{b2}. The COMMA benchmark, for instance, reveals weaknesses in state-of-the-art multimodal models, including proprietary models like GPT-4o, in collaborative tasks and agent-agent communication \cite{b1}. Similarly, the Multimodal Situational Safety (MSSBench) benchmark highlights current models' struggles with nuanced safety problems, emphasizing the need for more comprehensive evaluation frameworks \cite{b3}. The development of standardized benchmarks and evaluation methods is crucial for assessing MLLMs' performance and limitations, as well as identifying areas for future research and improvement.

Furthermore, biases and sensitivities in existing open-source models pose significant concerns for their reliability and trustworthiness in real-world applications \cite{b6}. The Multimodal Large Language Models to Support Real-World Fact-Checking survey proposes a framework for assessing MLLMs' capacity in fact-checking, revealing biases and weaknesses in current state-of-the-art models \cite{b6}. These findings highlight the importance of considering safety and security concerns in the development and deployment of MLLMs, as well as the need for more robust and transparent evaluation methods.

The technical details and methodologies employed in these studies also underscore the complexity of addressing data quality, bias, and security concerns in MLLM development. The COMMA benchmark, for example, uses a combination of open-source and closed-source models to evaluate collaborative performance in agent-agent and agent-human collaborations \cite{b1}. In contrast, the MSSBench dataset consists of 1,000 multimodal examples with annotated safety labels, using a combination of automated and human evaluation methods \cite{b3}. These diverse approaches highlight the need for a more holistic understanding of MLLM development, one that considers the interplay between models, data, and evaluation frameworks.

In conclusion, the challenges and limitations associated with data quality, bias, and security concerns in MLLM development are significant and multifaceted. Addressing these challenges will require a concerted effort to develop more comprehensive and standardized evaluation methods, as well as a greater emphasis on co-development strategies that prioritize high-quality data and transparent model development \cite{b4}, \cite{b5}. By acknowledging and addressing these limitations, researchers can work towards creating more reliable, trustworthy, and effective MLLMs that can be safely deployed in real-world applications. Ultimately, the successful integration of MLLMs and agents will depend on our ability to navigate these challenges and develop more robust, transparent, and accountable models that can support a wide range of collaborative tasks and applications \cite{b1}, \cite{b2}, \cite{b3}, \cite{b6}.

\section{Advanced Topics: Explainability, Transfer Learning, and Few-Shot Learning in Multimodal Models}
The realm of multimodal large language models and agents has witnessed significant advancements in recent years, with a growing emphasis on explainability, transfer learning, and few-shot learning. As highlighted by \cite{b1}, the foundations of multimodal machine learning are built upon key principles such as modality heterogeneity, connections, and interactions, which pose six core technical challenges including representation, alignment, reasoning, generation, transference, and quantification. The importance of incorporating multimodality to enhance Large Language Models' (LLMs) generation capabilities is further emphasized by \cite{b2}, who focus on methods that retrieve multimodal knowledge for augmented generation.

The intersection of LLMs with multimodal learning has led to significant advancements in generation and editing tasks across various domains, including image, video, 3D, and audio, as explored by \cite{b3}. Furthermore, the development of large language models extended into the multimodal domain has enabled AI agents to handle diverse user queries, as reviewed by \cite{b4}, who introduce essential components, categorize current research, and propose a framework for standardized evaluations. However, the complexity of these models also raises concerns about interpretability and explainability, which are addressed by \cite{b5} through a novel framework that categorizes existing research into data, model, and training \& inference perspectives.

A common theme across these studies is the integration of multimodal information to enhance AI capabilities, whether in generation, understanding, or interaction tasks. The need for explainability and interpretability is also a recurring theme, highlighting the importance of methodologies that improve transparency and trustworthiness in complex multimodal models. Moreover, the diversity of applications for multimodal AI is evident, ranging from visual question answering and cross-modal retrieval to image-text generation and large-scale agent interactions.

Despite these advancements, challenges persist, particularly in standardizing evaluation methods, which is crucial for comparing different models or approaches. As noted by \cite{b1}, the absence of standardized evaluation frameworks complicates progress in the field. Moreover, the integration of multiple modalities increases model complexity, posing challenges for scalability, interpretability, and computational efficiency. The explainability gap also remains a significant concern, as making multimodal models transparent and explainable is essential for trustworthiness and reliability, especially in high-stakes applications.

Methodological differences are also apparent, with some studies focusing on architectural innovations, such as CLIP/T5-based methods, while others emphasize training strategies or data-centric approaches. The scope of multimodality also varies, with some surveys focusing broadly on multimodal machine learning and others narrowing down to specific applications. Prioritization of explainability is another area of contrast, reflecting different priorities in the field.

Technical details and methodologies are also noteworthy, with various architectural designs for multimodal models, including transformers and their variants tailored for multimodal tasks. Training strategies such as multi-task learning, self-supervised learning, and the use of large-scale datasets to improve model generalizability are also explored. A range of evaluation metrics is proposed or discussed, reflecting the complexity of assessing performance in multimodal tasks.

In conclusion, the advanced topics of explainability, transfer learning, and few-shot learning in multimodal models are crucial for enhancing the capabilities of multimodal large language models and agents. While significant progress has been made, challenges persist, and addressing these challenges will be essential for realizing the full potential of multimodal AI systems. By highlighting key themes and developments, discussing implications and connections, and emphasizing the need for standardized evaluation frameworks and explainability, this analysis provides a comprehensive foundation for further research and development in this field, as envisioned by \cite{b1} and \cite{b5}. Ultimately, the cooperation between multimodal large language models and agents will depend on advancements in these areas, enabling more effective, transparent, and trustworthy interactions between humans and AI systems.

\section{Future Directions and Open Challenges in Multimodal Large Language Models}
The future directions and open challenges in multimodal large language models (MLLMs) are multifaceted and rapidly evolving, with significant implications for artificial general intelligence, world models, and human-computer interaction \cite{b1}. As highlighted by recent surveys and reviews \cite{b2}, \cite{b3}, a key challenge in this field is the development of standardized evaluation methodologies and benchmarks, such as COMMA \cite{b4}, to assess the performance of MLLMs in collaborative tasks and communicative settings. This need for standardization is critical because it will enable meaningful comparisons between different models and facilitate the identification of areas where improvements are needed.

Another crucial area of research is improving collaboration between agents and humans through more effective communication and cooperation frameworks \cite{b5}. Techniques such as Multimodal Chain of Thought (M-COT) \cite{b6}, Multimodal Instruction Tuning (M-IT) \cite{b7}, and Multimodal In-Context Learning (M-ICL) \cite{b8} have shown promise in enhancing the performance of MLLMs in multimodal tasks. Furthermore, incorporating external rule systems and embodied intelligence can significantly enhance world simulation capabilities and reasoning in MLLMs \cite{b9}. However, addressing limitations in decision-making and improving the ability of agents to understand and respond appropriately to human input remain open challenges \cite{b10}.

The exploration of new techniques and architectures for multimodal learning is also a fertile ground for future research. For instance, extending techniques like LAVR (Learning to Abstract and Verify Representations) \cite{b11} could provide more robust representations across different modalities. Additionally, investigating the applications of MLLMs in various domains, including image, video, 3D, and audio generation and editing, as well as their potential role in artificial general intelligence \cite{b12}, will be crucial for unlocking their full potential.

A common theme across these future directions is the emphasis on multimodal learning and its potential to enhance AI capabilities \cite{b13}. The ability of MLLMs to integrate information from multiple sources and modalities makes them particularly suited to tasks that require a deep understanding of context and the ability to reason abstractly. However, achieving true multimodal understanding will require overcoming significant technical hurdles, including the development of unified models that can seamlessly integrate different modalities \cite{b14}.

In conclusion, the future of MLLMs is marked by both significant promise and substantial challenges. Addressing these challenges through continued research into standardized evaluation methodologies, improved collaboration between agents and humans, enhanced reasoning capabilities, and innovative architectures for multimodal learning will be essential for realizing the potential of MLLMs \cite{b15}. As researchers continue to push the boundaries of what is possible with MLLMs, it is likely that we will see major advancements in areas ranging from human-computer interaction to artificial general intelligence, ultimately leading to more sophisticated and capable AI systems.

\section{Conclusion and Emerging Trends in Multimodal AI Research}
The realm of multimodal AI, particularly focusing on how large language models (LLMs) and agents can cooperate, is rapidly evolving \cite{b1}. Recent surveys and research papers have delved into various aspects of this cooperation, including multimodal generation, editing, and the development of benchmarks for evaluating multimodal multi-agent systems \cite{b2}. A key finding in this area is the emergence of Multimodal Large Language Models (MLLMs) as a new research hotspot, with capabilities such as writing stories based on images and OCR-free math reasoning, indicating a potential path towards artificial general intelligence \cite{b3}. The introduction of benchmarks like COMMA (Communicative Multimodal Multi-Agent Benchmark) emphasizes the importance of language-based communication between agents in collaborative tasks, revealing surprising weaknesses in state-of-the-art models when it comes to agent-agent collaboration \cite{b4}.

A common theme across papers is the emphasis on integrating different modalities and enabling cooperation between agents or models to achieve more complex tasks \cite{b5}. There's a recognized need for benchmarks that can effectively evaluate the performance of multimodal systems, especially in multi-agent scenarios \cite{b6}. Papers frequently highlight challenges related to modality heterogeneity, alignment, and interaction, underscoring the complexity of dealing with different data types and sources \cite{b7}. The debate on the focus between developing models for specific tasks (like image-based story writing) versus aiming for more generalized artificial intelligence that can handle a wide range of multimodal inputs and tasks is also a notable aspect of current research \cite{b8}.

Technical details, such as MLLM architectures, training strategies, and data requirements, are crucial in advancing the field \cite{b9}. Innovations like Multimodal ICL (M-ICL), Multimodal CoT (M-CoT), and LLM-Aided Visual Reasoning (LAVR) have been highlighted as significant advancements \cite{b10}. Furthermore, there's an emphasis on developing appropriate evaluation metrics for multimodal models, considering aspects beyond traditional accuracy measures to include generative capabilities, safety, and ethical considerations \cite{b11}.

Despite the progress made, challenges related to data heterogeneity and availability remain significant, with a need for large, diverse datasets for training MLLMs \cite{b12}. The complexity of multimodal models can also make them less explainable, posing challenges for understanding decision-making processes \cite{b13}. Ensuring the safety and ethical use of these models becomes a pressing concern as generative capabilities improve, requiring careful consideration of potential biases and misuse scenarios \cite{b14}.

In conclusion, the cooperation between LLMs and agents in multimodal AI is an area of vibrant research activity. While significant progress has been made, challenges necessitate continued innovation and investment in this field. Future research directions should prioritize the development of more sophisticated benchmarks, addressing the technical challenges outlined in recent taxonomies, and ensuring that advancements are aligned with ethical standards and safety protocols \cite{b15}. The integration of multimodal AI with other emerging technologies is expected to have profound implications for various applications and domains, underscoring the need for multidisciplinary research efforts to fully harness its potential \cite{b16}. Ultimately, the goal of achieving more generalized and explainable multimodal AI systems that can effectively cooperate with humans and other agents will require sustained research efforts and collaborations across academia, industry, and policy-making communities \cite{b17}.

\section{Technical Appendix: Detailed Model Architectures and Algorithms}
The technical aspects of multimodal large language models and agent cooperation have been extensively explored in recent surveys, providing valuable insights into the development of complex AI systems \cite{b1}. A key finding in the survey paper "Large Multimodal Agents: A Survey" is the importance of developing large multimodal agents (LMAs) that can interpret and respond to diverse multimodal user queries \cite{b2}. This concept is further reinforced by the paper "A Survey on Game Playing Agents and Large Models", which emphasizes the need for collaborative frameworks that integrate multiple LMAs to enhance collective efficacy \cite{b3}. The introduction of the COMMA benchmark in the paper "COMMA: A Communicative Multimodal Multi-Agent Benchmark" provides a comprehensive evaluation framework for multimodal multi-agent systems through language communication, highlighting the significance of standardized assessment methodologies \cite{b4}.

A common theme across the surveyed papers is the importance of multimodal interaction, including text, images, video, and audio, in developing complex AI systems \cite{b5}. The need for effective collaboration and communication between agents, as well as between agents and humans, is a recurring pattern in the papers, underscoring the significance of designing frameworks that facilitate seamless interaction \cite{b6}. Furthermore, the surveys highlight the diversity of evaluation methodologies used across existing studies, emphasizing the need for standardized frameworks to facilitate meaningful comparisons and inform future research directions \cite{b7}.

The debate between centralized and decentralized architectures is a notable aspect of the surveyed papers, with some advocating for centralized approaches due to their potential for enhanced control, while others propose decentralized methods, which offer improved scalability \cite{b8}. Additionally, the surveys touch on the contrast between symbolic and connectionist AI, with some papers emphasizing the importance of symbolic representations in multimodal interaction, while others highlight the benefits of connectionist approaches \cite{b9}. The application of large language models in agent-based modeling and simulation, as explored in the survey paper "Large Language Models Empowered Agent-based Modeling and Simulation", has significant implications for enhancing simulation capabilities and provides a promising direction for future research \cite{b10}.

In terms of technical details, the papers discuss various model architectures, including transformer-based models, convolutional neural networks (CNNs), and recurrent neural networks (RNNs) \cite{b11}. The surveys also cover different training methods, such as self-supervised learning, supervised learning, and reinforcement learning, highlighting the importance of selecting appropriate training approaches to achieve optimal performance \cite{b12}. Cross-modal learning is another critical aspect, with the papers emphasizing the challenges of effective feature fusion and alignment across different modalities \cite{b13}.

Despite the advancements in multimodal large language models and agent cooperation, several limitations and challenges remain. Scalability and robustness are significant concerns, as existing frameworks often struggle to handle complex, real-world scenarios \cite{b14}. The development of standardized evaluation metrics is also crucial, as current methodologies often lack consistency, making it difficult to compare the performance of different LMAs and multimodal multi-agent systems \cite{b15}. Furthermore, explainability and transparency are essential aspects of multimodal AI systems, with techniques needed to provide insights into decision-making processes and enhance trust in these complex systems \cite{b16}. By addressing these challenges and building on existing research, future studies can develop more effective and efficient multimodal large language models and agent cooperation frameworks, ultimately leading to significant advancements in the field of artificial intelligence.

\section{References and Bibliography of Key Works in the Field}
The field of multimodal large language models and agent cooperation has witnessed significant advancements, with various studies exploring their applications, methodologies, and challenges. A survey of key papers in this field reveals a plethora of research focusing on large multimodal agents, communicative multimodal multi-agent benchmarks, game playing agents, agent-based modeling and simulation, and multimodal generation and editing. For instance, the concept of large multimodal agents (LMAs) is introduced in \cite{b1}, providing a systematic review of LLM-driven multimodal agents, which categorizes current research into four distinct types and reviews collaborative frameworks, highlighting the need for standardized evaluation methodologies.

The introduction of the COMMA benchmark \cite{b2} to evaluate the collaborative performance of multimodal multi-agent systems through language communication has revealed surprising weaknesses in state-of-the-art models, especially in agent-agent collaboration scenarios. Furthermore, a survey on game playing agents and large models \cite{b3} explores their capabilities and potential in complex game scenarios, identifying challenges and future research avenues for the advancement of LMs in games. The integration of large language models into agent-based modeling and simulation is examined in \cite{b4}, discussing challenges such as environment perception, human alignment, action generation, and evaluation, and providing an overview of recent works across four domains: cyber, physical, social, and hybrid.

A common theme across these papers is the integration of large language models with various applications, including multimodal tasks, agent-based modeling, game playing, and generation/editing tasks \cite{b5}. The need for standardized evaluation methodologies to effectively compare different models and approaches is a recurring mention \cite{b6}. Challenges such as understanding, generating, and editing multimodal content are highlighted across several papers, emphasizing the complexity of these tasks \cite{b7}. The potential applications of large multimodal agents and related technologies in real-world scenarios, including games, simulations, and human-computer interaction, are frequently discussed \cite{b8}.

While some papers focus primarily on the understanding aspect of multimodal tasks, others delve into generation and editing, showing a contrast in approach towards leveraging large language models \cite{b9}. The introduction of specific benchmarks like COMMA indicates a shift towards more structured evaluation methods, contrasting with earlier, more varied approaches to assessing model performance \cite{b10}. Papers vary in their focus on specific domains (e.g., games, simulations), reflecting different viewpoints on where large multimodal agents can provide the most value \cite{b11}.

The technical details of collaborative frameworks for large multimodal agents are explored, including how these frameworks facilitate cooperation between agents \cite{b12}. Various papers discuss the use of LLMs like CLIP and T5 for multimodal tasks, highlighting their architectures and how they are adapted for generation and editing tasks \cite{b13}. The technical aspects of generative models used in multimodal generation and editing, such as GANs and VAEs, are discussed, along with their integration with LLMs \cite{b14}.

However, several limitations and challenges are noted, including the difficulty in evaluating the performance of large multimodal agents due to the lack of standardized metrics \cite{b15}. The inherent complexity of understanding and generating multimodal content poses significant technical challenges \cite{b16}. Generative AI safety is identified as a critical area of concern, particularly with the potential for misuse of generated content \cite{b17}. As models grow in size and complexity, scalability and efficiency become significant challenges that need to be addressed through innovative architectural designs and computational strategies \cite{b18}.

In conclusion, this analysis provides a comprehensive overview of the current state of research in large multimodal agents and their applications, highlighting key findings, common themes, contrasting viewpoints, technical details, and limitations. It serves as a foundation for further exploration into the potential of integrating large language models with multimodal tasks and agent cooperation, aiming to advance the field towards more sophisticated and applicable technologies \cite{b19}. The implications of this research are far-reaching, with potential applications in various domains, and it is essential to continue addressing the challenges and limitations to fully realize the potential of large multimodal agents and their cooperative capabilities \cite{b20}.

\section{Conclusion}
In conclusion, the analysis of papers related to multimodal large language models and agent cooperation reveals a complex and multifaceted landscape, with significant progress made in recent years \cite{b1}. The importance of communication and collaboration between agents in achieving tasks beyond individual capabilities is a recurring theme, underscoring the need for effective coordination and interaction mechanisms \cite{b2}. The development of large multimodal agents (LMAs) has been a key area of research, with these models demonstrating superior performance in powering text-based AI agents and interpreting diverse multimodal user queries \cite{b3}. Furthermore, LLM-based multi-agent reinforcement learning frameworks have shown great potential, although challenges persist in addressing coordination and communication between agents \cite{b4}.

The symmetry of agents and interplay with prompts has also been identified as a critical factor in unlocking the reasoning capability of large language models, with empirical results demonstrating that carefully developed prompt engineering can approach state-of-the-art performance of complex multi-agent mechanisms \cite{b5}. Additionally, multimodal generation and editing have seen notable advancements, with LLMs playing a crucial role in various domains, including image, video, 3D, and audio \cite{b6}. The COMMA benchmark has highlighted surprising weaknesses in state-of-the-art models, including proprietary models like GPT-4o, in agent-agent collaboration tasks, emphasizing the need for more effective evaluation frameworks \cite{b7}.

Despite these advancements, significant challenges remain, including the lack of standardization in evaluation methodologies and the limited availability of datasets and technical components for multimodal generation and editing \cite{b8}. Moreover, existing agentic benchmarks fail to address key aspects of inter-agent communication and collaboration, and state-of-the-art models often struggle to outperform simple random agent baselines in agent-agent collaboration tasks \cite{b9}. To overcome these limitations, future research should focus on developing more comprehensive evaluation frameworks and addressing the technical challenges associated with multimodal large language models and agent cooperation.

The implications of this research are far-reaching, with potential applications in areas such as human-computer interaction, natural language processing, and artificial intelligence \cite{b10}. As the field continues to evolve, it is essential to address the limitations and challenges identified in this analysis, including the need for more effective coordination and communication mechanisms, standardized evaluation frameworks, and improved multimodal generation and editing capabilities. By doing so, researchers can unlock the full potential of multimodal large language models and agent cooperation, enabling the development of more sophisticated and effective AI systems \cite{b11}. Ultimately, this research has the potential to transform the way we interact with technology, enabling more natural, intuitive, and effective human-computer interactions \cite{b12}.

\end{document}